% SOURCE_AUTHORITY: sga5_fr_workpass.tex, lines 13477--13519 (LF-normalized unit).
% SOURCE_SHA256: 791F4EFFC5E02832D5D77ED03518C8156D6F07E4C8238B03545DB93D883FBB28
La demostración de la proposición 4.6 utiliza esencialmente el lema siguiente:

\medskip
\noindent
\textbf{Lema 4.7.} \emph{Con las hipótesis y notaciones de la proposición 4.6, supóngase además que $\varphi = \id_G$, es decir, que la acción de $G$ sobre $R$ conmuta con $f'_R$, y que $r = \nu(f'_R) \geq 2$.}

\emph{Sea
\[
c = \card(\im(G \longrightarrow \Aut_k(\fm/\fm^2))) \quad .
\]
Se tiene entonces la relación de divisibilidad siguiente:
\[
c \mid (r - 1) \quad .
\]}

\medskip

\begin{prooftext}[Demostración]
Es claro que el entero $r$ satisface las condiciones siguientes:
\begin{enumerate}
\item[a)] el homomorfismo
\[
f'_r : A/\fm^r \longrightarrow A/\fm^r
\]
inducido por $f'_R$ es la identidad.
\item[b)] el homomorfismo
\[
f'_{r+1} : A/\fm^{r+1} \longrightarrow A/\fm^{r+1}
\]
es distinto de la identidad.
\end{enumerate}

Por otra parte, es evidente que $G$ opera sobre todos los espacios vectoriales $\fm^i/\fm^{i+1}$ sobre $k$.

Mostraremos primero que, bajo las hipótesis del lema, existe un único homomorfismo de $k[G]$-módulos
\[
u(f'_R) : \fm/\fm^2 \longrightarrow \fm^r/\fm^{r+1}
\]
tal que se satisfacen las dos condiciones siguientes:
\begin{enumerate}
\item[1)] $u(f'_R) \neq 0$
\item[2)] $u(f'_R)(x \bmod \fm^2) = (f'_R(x) - x) \bmod \fm^{r+1}$, cualquiera que sea $x \in \fm$.
\end{enumerate}
